\begin{abstract}
Latent heat thermal energy storage using phase change materials (PCMs) is an effective strategy for mitigating the temporal mismatch between energy supply and demand, but the inherently low thermal conductivity of paraffin-based PCMs severely limits their charging and discharging rates. In this study, we experimentally investigate the thermal performance of paraffin-based PCM (phase change temperature 42\textdegree C) infiltrated into additively manufactured AlSiMg triply periodic minimal surface (TPMS) Gyroid lattice structures. A nine-sensor array organized into three radial groups was used to characterize the spatial temperature field under three heating power levels (10~W, 20~W, and 30~W). The results reveal three distinct thermal regimes---sensible heating, phase change plateau near 42\textdegree C, and post-melting sensible heating---with the plateau duration decreasing from 5--8~minutes at 10~W to near-absence at 30~W. A consistent radial temperature stratification ($T_1 > T_{2\text{-}5} > T_{6\text{-}9}$) was observed, with the temperature spread between innermost and outermost sensor groups limited to 17.1\textdegree C at 10~W, demonstrating effective thermal conductivity enhancement by the TPMS lattice. The heating duration decreased from 25.4~min at 10~W to 8.5~min at 30~W. Systematic outlier analysis revealed position-dependent thermal anomalies in the outer sensor group. These results establish the Gyroid baseline for a comparative study of TPMS topologies (Gyroid, Primitive, and IWP), providing a foundation for the rational design of TPMS-enhanced PCM thermal energy storage systems.

\noindent\textbf{Keywords:} phase change material; triply periodic minimal surface; Gyroid; additive manufacturing; thermal energy storage; latent heat
\end{abstract}
